
\section{El Lema de Morse}

Comenzamos introduciendo los puntos críticos de una función suave y el concepto de índice de un punto crítico no degenerado. Concluiremos la sección con el primer pilar de la teoría de Morse y sobre lo que se edificarán gran parte de los resultados posteriores: el lema de Morse.

\begin{definition}
    Sea $f: \mathrm{M} \to \mathbb{R}$ una función suave.
    \begin{enumerate}[label=\roman*)]
        \item Diremos que $p\in \mathrm{M}$ es un \textit{punto crítico de $f$} si $d_{p}f:T_{p}M\to \mathbb{R}$ es idénticamente cero. Es decir, dado un sistema de coordenadas local $(x_{1}, \dots ,x_{n})$ en un entorno $\mathcal{U}$ de $p$ tenemos
        \begin{equation*}
            \frac{\partial f}{\partial x_1}(p)= \dots = \frac{\partial f}{\partial x_n}(p) = 0.
        \end{equation*}
        
        \item A su vez, diremos que $f(p)\in \mathbb{R}$ es un \textit{valor crítico de $f$} si $p \in \mathrm{M}$ es un punto crítico.
        
    \end{enumerate}
\end{definition}
\bigskip

\begin{definition}
    Sea $p\in \mathrm{M}$ un punto crítico de f , diremos que $p$ es \textit{no degenerado} si la matriz  $\left(\frac{\partial^{2}f}{\partial x_i \partial x_j}(p)\right)$ es no singular, es decir, su determinante es no nulo.
\end{definition}
\bigskip

\begin{definition}
    Sean $f:\mathrm{M}\to \mathbb{R}$ una función suave y $p\in M$ un punto crítico de f. Definimos $\textrm{Hess}_{p}f: T_{p}\mathrm{M}\to \mathbb{R}$ como la forma bilineal simétrica sobre el plano tangente a $\mathrm{M}$ en $p$ construida como sigue:
    \begin{quote}
        Tomamos $v,w\in T_{p}\mathrm{M}$ y sean $\tilde{v}$ y $\tilde{w}$ extensiones locales suyas a campos vectoriales. Entonces 
        \begin{equation*}
            \text{Hess}_{p}f(v,w)=\tilde{v}_{p}(\tilde{w}_{p}(f)),
        \end{equation*}
        donde $\tilde{v}_{p}$ es simplemente $v$ y $\tilde{w}_{p}(f)$ denota la derivada parcial de $f$ en la dirección de $\tilde{w}$.
    \end{quote}
\end{definition}
\bigskip
\begin{remark}
    Es rápido comprobar que $\text{Hess}_{p}f$ es en efecto simétrica y está bien definida:

    \begin{enumerate}[label=(\roman*)]
        \item Es simétrica ya que $\tilde{v}_{p}(\tilde{w}_{p}(f)) - \tilde{w}_{p}(\tilde{v}_{p}(f)) = [\tilde{v},\tilde{w}]_{p}(f) = 0$ donde $[\tilde{v},\tilde{w}]$ es el corchete de Lie de $\tilde{v}$ y $\tilde{w}$ cuya expresión en coordenadas recordemos que es
        \begin{equation*}
            [\tilde{v},\tilde{w}]_{p}(f)=\sum_{i=1}^n \big( \tilde{v}_p(\tilde{w}_{i}) - \tilde{w}_p(\tilde{v}_{i}) \big) \left. \frac{\partial}{\partial x_{i}}\right|_p f.
        \end{equation*}
        Luego se anula al ser $p$ un punto crítico.

        \item Por otro lado, está bien definida ya que $\tilde{v}_{p}(\tilde{w}_{p}(f))$ no depende de las extensiones de $v$ y $w$ elegidas puesto que nos importa el vector del campo en el punto $p$ y este no cambia.
    \end{enumerate}
\end{remark}
\bigskip

Veamos ahora como es $\text{Hess}_pf$ en coordenadas. De nuevo, tomemos un sistema de coordenadas local $(x_{1}, \dots ,x_{n})$, tenemos
\begin{equation*}
    v=\sum_{i=0}^n a_i \left.\frac{\partial}{\partial x_i}\right|_p ~~\text{y}~~ w=\sum_{i=0}^n b_i \left.\frac{\partial}{\partial x_i}\right|_p .
\end{equation*}

Y para las extensiones $\tilde{v}$ y $\tilde{w}$ basta tomar como coeficientes las funciones constantes $a_i$ y $b_j$. Así
\begin{equation*}
    \text{Hess}_pf= \tilde{v}_{p}(\tilde{w}_{p}(f)) = v\left(\sum_{i=0}^n b_i \left.\frac{\partial}{\partial x_i}\right|_p f\right)=\sum_{i,j} a_ib_i\left.\frac{\partial^{2}}{\partial x_i \partial x_j}\right|_p f .
\end{equation*}

Luego la matriz $\left(\frac{\partial^{2}f}{\partial x_i \partial x_j}(p)\right)$ mencionada antes representa la forma bilineal $\text{Hess}_pf$ respecto de la base $\left( \left.\frac{\partial}{\partial x_i}\right|_p,\dots, \left.\frac{\partial}{\partial x_n}\right|_p \right)$ de $T_p\mathrm{M}$.
\bigskip

\begin{definition}
    Sean $V$ un espacio vectorial y $H$ una forma bilineal sobre $V$. Entonces
    
    \begin{enumerate}[label=\roman*)]
        \item El \textit{índice} de $H$ es la máxima dimensión de un subespacio de $V$ sobre el cual $H$ es definida negativa.
        
        \item La \textit{nulidad} de $H$ es la dimensión del radical de $H$, es decir,
        \begin{equation*}
            \text{dim}(\text{Rad}(H)) ~\text{ con }~ \text{Rad}(H)=\{v\in V:H(v,w)=0, ~\forall w\in V \}.
        \end{equation*}
        
    \end{enumerate}
\end{definition}

\begin{remark}
    Un punto crítico $p\in M$ de $f$ suave es no degenerado si y solo si la nulidad de $H$ sobre $T_pM$ es $0$. En efecto:
    \vspace{-2em}
    \begin{quote}
        \mbox{}\\[0.5em]
        $\boxed{\implies}$ Supongamos que existe $v\in \text{Rad}(H)$ distinto del cero, entonces $H(v,w)=v^tHw=0$ para todo $w\in T_pM$, luego $v^tH=0 \text{ si y solo si } v^t = 0$ ya que $H$ es no singular y $v^tH$ es una combinación lineal las filas de $H$, contradicción.

        $\boxed{\impliedby}$ $\text{Rad}(H)=\{0\}$. Supongamos que $p$ es degenerado y sea $v\in T_pM\setminus\{0\}$, entonces existe $w\in T_pM$ tal que $v^tHw\neq 0$. Así, $v^tH\neq 0$ y como $v$ era arbitrario se tiene que $v^tH\neq 0$ para todo $v\in T_pM\setminus\{0\}$, luego $H$ es no singular y hemos llegado a una contradicción.
    \end{quote}
\end{remark}
\bigskip

\begin{nota}
    De ahora en adelante nos referiremos a la forma bilineal inducida por $f$ como $\text{Hess}_pf$ y a su matriz con respecto de la base del plano de tangente como $H_pf$ o simplemente $H$ si el contexto lo permite. Además nos referiremos al índice y la nulidad de $\text{Hess}_pf$ como \textit{índice y nulidad de $f$ en $p$}.
\end{nota}
\bigskip


Ahora que hemos desarrollado el vocabulario para hablar de los puntos críticos de una función sobre los reales, cerramos esta sección con el enunciado y la demostración del \textit{Lema de Morse}. 

Este resultado es fundamental para el desarrollo de la teoría de Morse. A grandes rasgos, establece que en un entorno de un punto crítico no degenerado existe un sistema local de coordenadas respecto al cual la función $f$ puede expresarse como una forma cuadrática. Desde un punto de vista geométrico, esto implica que $f$ se asemeja a un paraboloide en dicho entorno. Esto resulta increíblemente útil y constituye la base de los argumentos topológicos que usaremos a partir de este punto.

\bigskip

Antes de enunciar el Lema de Morse necesitamos probar el siguiente resultado de cálculo diferencial.

\begin{lemma}
    Sean $V$ un entorno convexo del origen en $\mathbb{R}^n$ y $f\in \mathcal{C}^\infty(V)$ con $f(0)=0$. Entonces,
    \begin{equation*}
        f(x_1,\dots,x_n)=\sum_{i=1}^n x_ig_i(x_1,\dots,x_n)
    \end{equation*}
    para unas funciones $g_i\in \mathcal{C}^\infty(V)$ que además, cumplen que $g_i(0)=\frac{\partial f}{\partial x_i}(0)$.
\end{lemma}
\begin{proof}
    Se tiene
\begin{align*}
    f(x_i,\dots,x_n)=\int_0^1 \frac{df(tx_1,\dots,tx_n)}{dt}dt = \int_0^1 \big(\sum_{i=1}^n \frac{\partial f}{\partial x_i}&(tx_1,\dots,tx_n)x_i \big)dt =\\
    &=\sum_{i=1}^n x_i \left(\int_0^1 \frac{\partial f}{\partial x_i}(tx_1,\dots,tx_n )dt\right)
\end{align*}

Tomamos $g_i(x_1,\dots,x_n) = \int_0^1 \frac{\partial f}{\partial x_i}(tx_1,\dots,tx_n )dt$ y se cumple que $g_i(0)=\int_0^1 \frac{\partial f}{\partial x_i}(0)dt =\frac{\partial f}{\partial x_i}(0)$.

\end{proof}
\bigskip

Vayamos ahora con la demostración, para la que nos hemos basado tanto en la de Milnor (\cite[p.~6]{milnor1963}) como en la que se encuentra en \cite[p.~431]{marsdenHoffman1993eca}.

\begin{lemma}[Lema de Morse]
\label{lemamorse}
    Sean $p$ un punto crítico no degenerado de una función suave $f:M\to \mathbb{R}$. Entonces existe un sistema de coordenadas local $(y_1,\dots, y_n)$ en un entorno $\mathcal{U}$ de $p$ tal que $y_i(p)=0, \forall i\in \{1\dots,n\}$ y se verifica
    \begin{equation*}
        f=f(p)-(y_1)^2-\dots-(y_\lambda)^2 + (y_{\lambda +1})^2+\dots + (y_n)^2
    \end{equation*}
    en todo $\mathcal{U}$ con $\lambda$ el índice de $f$ en $p$.
\end{lemma}
\bigskip

\begin{proof}[Demostración del Lema de Morse]
    %################ PARTE 1, INDICE ######################
    En primer lugar, veamos que de poder expresar $f$ de esa manera en $\mathcal{U}$ entonces $\lambda$ es necesariamente el índice de $f$ en $p$. Sea $(x_1,\dots, x_n)$ un sistema de coordenadas local en $\mathcal{U}$ y supongamos que
    \begin{equation*}
        f=f(p)-(x_i)^2-\dots-(x_\lambda)^2 + (x_{\lambda +1})^2+\dots + (x_n)^2.
    \end{equation*}
    Entonces se tiene que 

    \[
        \frac{\partial^2 f}{\partial x_i \partial x_j}=
        \begin{cases}
            -2, ~i=j\le \lambda\\
            ~~~2, ~i=j\ge \lambda\\
            ~~~0, i\neq j
        \end{cases}
    \]

    Es decir, $H_pf$ es una matriz diagonal con $\lambda$ valores negativos y $n-\lambda$ valores positivos. Luego tiene $\lambda$ autovalores negativos, lo que quiere decir que existe un subespacio de dimensión $\lambda$ de $T_pM$ en el que es definida negativa y otro de dimensión $n-\lambda$ en el que es definida positiva.

    La única forma en que $\lambda$ pudiera ser distinto que el índice de $f$ en $p$ es si existiera un subespacio de dimensión mayor que el mencionado en el que $H_pf$ fuese definida negativa, lo cual es imposible.\\


    %############ PARTE 2, SISTEMA DE COORDENADAS ##################
    Hacen falta algunos preparativos antes de demostrar que el sistema de coordenadas $(y_1,\dots,y_n)$ existe.
    
    En primer lugar podemos asumir sin pérdida de generalidad que, en coordenadas, $p$ es el origen de $\mathbb{R}^n$ componiendo la carta escogida con la traslación que lleva $p$ al origen, que es un difeomorfismo. También asumimos $f(p)=f(0)=0$, de no ser así podríamos tomar $\bar{f}(q)=f(q)-f(p)$ que verifica $\bar{f}(p)=0$ y tiene a $p$ como punto crítico no degenerado y deshacer después el cambio. 

    Por el lema anterior $f(x_1,\dots,x_n)=\sum_{j=1}^n x_jg_j(x_1,\dots,x_n)$ en un entorno convexo del origen. Ahora, como $0$ es un punto crítico de $f$ se tiene
    \begin{equation*}
        \frac{\partial f}{\partial x_j}(0)=0 ~\text{ y }~ g_j(0)=\int_0^1 \frac{\partial f}{\partial x_j}(0)~dt=0.
    \end{equation*}
    Luego podemos aplicar de nuevo el lema a cada una de las $g_j$ obteniendo
    \begin{equation*}
        g_j(x_1,\dots,x_n)=\sum_{i=1}^n x_ih_{ij}(x_1,\dots,x_n).
    \end{equation*}
    Y, por tanto,
    \begin{equation*}
        f(x_1,\dots,x_n)=\sum_{i,j} x_ix_jh_{ij}(x_1,\dots,x_n).
    \end{equation*}
    %aqui iba lo de quotes
    Además, la matriz $H(0)=(h_{ij}(0))$ es igual a $\frac{1}{2}H_0f=\frac{1}{2}(\frac{\partial^2f}{\partial x_i\partial x_j}(0))$ y por lo tanto es simétrica y no singular. En efecto, sabemos que 
    \begin{equation*}
        h_{ij}=\int_0^1 \frac{\partial g_i}{\partial x_j}(tx_1,\dots,tx_n)dt ~~\text{ y}~~ g_{i}=\int_0^1 \frac{\partial f}{\partial x_i}(tx_1,\dots,tx_n)dt.
    \end{equation*}
    Luego
    \begin{gather*}
        \frac{\partial g_i}{\partial x_j}(x_1,\dots,x_n)= \int_0^1 \frac{\partial}{\partial x_j}\left(\frac{\partial f}{\partial x_i}(tx_1,\dots,tx_n)\right)dt = \int_0^1 \left(\sum_{k=1}^n\frac{\partial^2 f}{\partial x_i \partial x_k}(tx_1,\dots,tx_n) \frac{\partial(tx_k)}{\partial x_j}\right)dt =\\
        = \int_0^1t\frac{\partial^2 f}{\partial x_i \partial x_j}(tx_1,\dots,tx_n)dt.
    \end{gather*}
    Así, 
    \begin{gather*}
        \frac{\partial g_i}{\partial x_j}(0)=\int_0^1 t\frac{\partial^2 f}{\partial x_i \partial x_j}(0)dt=\frac{\partial^2 f}{\partial x_i \partial x_j}(0)\int_0^1tdt= \frac{1}{2}\frac{\partial^2 f}{\partial x_i \partial x_j}(0)\implies\\
        \implies h_{ij}(0)=\frac{1}{2}\frac{\partial^2 f}{\partial x_i \partial x_j}(0)\int_0^1 dt =\frac{1}{2}\frac{\partial^2 f}{\partial x_i \partial x_j}(0)
    \end{gather*}
    %############### INDUCCIÓN ###########################
    Finalmente, tenemos $f$ expresada similarmente a una forma cuadrática. Lo que queremos hacer a continuación es diagonalizarla, para ello procederemos por inducción. 
    
    El caso base de la inducción ($r=1$) ya está demostrado.

    Supongamos ahora que existe un sistema de coordenadas $(u_i,\dots,u_n)$ en un entorno $\mathcal{U}_1$ de $0$ quizás más pequeño que $\mathcal{U}$ tal que
    \begin{equation}
        f = \pm (u_1)^2 \pm\dots\pm (u_{r-1})^2 + \sum_{i,j\ge r} u_i u_j H_{ij}(u_1,\dots,u_n)
    \label{eq1}
    \end{equation}
    en todo $\mathcal{U}_1$ con $r > 1$ y $H_{ij}=H_{ji}$ para cualesquiera $i$ y $j$.

    Podemos realizar ahora un cambio lineal de coordenadas en $(u_r,\dots, u_n)$ para diagonalizar $\sum_{i,j\ge r} u_i u_j H_{ij}(0)$.
    
    Sabemos que, por ser la matriz $(H_{ij})$ no singular, los términos de la diagonal son todos no nulos. Luego podemos asumir $H_{rr}\neq 0$. Sea $g(u_1,\dots,u_n)=|H_{rr}(u_1,\dots,u_n)|^{1/2}$, es claro que existe un entorno $\mathcal{U}_2\subset \mathcal{U}_1$ de $0$ en el que es suave y no se anula. Definimos el nuevo sistema de coordenadas
    \begin{equation}
        \begin{cases}
            V_i(u_1,\dots,u_n) = u_i,~~i\neq r\\
            V_r(u_1,\dots,u_n)=g(u_1,\dots,u_n)\left[u_r+\sum_{i>r}u_i\frac{H_{ir}(u_1,\dots,u_n)}{H_{rr}(u_1,\dots,u_n)}\right],
        \end{cases}
    \label{eq2}
    \end{equation}

    veamos que en efecto lo es. Se tiene que la matriz Jacobiana de la transformación de coordenadas en el $0$ es
    \begin{equation*}
        \frac{\partial (V_1,\dots, V_n)}{\partial (u_1,\dots, u_n)} = 
        \begin{pmatrix}
            1 & 0 & ~ & \dots & ~ & ~ & 0 \\
            0 & ~ & ~ & ~ & ~ & ~ & 0 \\
            \vdots & ~ & 1 & ~ & ~ & ~ & \vdots \\
            \frac{\partial V_r}{\partial u_1} & \dots & ~ & g(0) & ~ & \dots & \frac{\partial V_r}{\partial u_n} \\
            0 & ~ & ~ & ~ & 1 & ~ & 0 \\
            \vdots & ~ & ~ & ~ & ~ & ~ & \vdots \\
            0 & ~ & ~ & \dots & ~ & 0 & 1
        \end{pmatrix},
    \end{equation*}
    y es no singular, basta restar a la $r$-ésima fila las demás multiplicadas por el término $\frac{\partial V_i}{\partial u_j}$ adecuado para que sea diagonal, y los términos de la diagonal son todos no nulos.

    Se sigue entonces del \textit{teorema de la función inversa} que la transformación es un difeomorfismo sobre un entorno más pequeño $\mathcal{U}_3\subset \mathcal{U}_2$ del $0$ y por tanto $(v_1,\dots, v_n)$ es un sistema de coordenadas local alrededor de $p$.

    Ahora consideramos
    \begin{equation}
        u_ru_rH_{rr}(u_1,\dots, u_n) + 2\sum_{j=r+1}^n u_ju_rH_{jr}(u_1,\dots, u_n).
    \label{eq3}
    \end{equation}

    Usando la simetría de $H_{ij}$ observamos esta expresión no es más que los términos de \eqref{eq1} que tienen $j=r$ o $i=r$. Comparando \eqref{eq1} y \eqref{eq2} llegamos a que \eqref{eq3} es
    \begin{equation}
        \pm V_rV_r - \frac{1}{H_{rr}}\left[ \sum_{i>r} u_iH_{ir}(u_1,\dots, u_n) \right]^2
    \label{eq4}
    \end{equation}

    En efecto,
    \begin{align*}
        V_rV_r =& ~ g^2 \left[u_r+\sum_{i>r}u_i\frac{H_{ir}}{H_{rr}}\right]^2 =\\
        =& ~ g^2 \left( u_r^2 + \left[ \sum_{i>r}u_i\frac{H_{ir}}{H_{rr}}\right]^2 + 2u_r\sum_{i>r}u_i\frac{H_{ir}}{H_{rr}} \right) = \\
        =& ~ |H_{rr}| u_r^2 + \frac{|H_{rr}|}{H_{rr}^2} \left[\sum_{i>r}u_iH_{ir}\right]^2 + 2u_r\frac{|H_{rr}|}{H_{rr}}\sum_{i>r}u_iH_{ir} = \\
        =& ~ \pm H_{rr}u_r^2 \pm \frac{1}{H_{rr}}\left[\sum_{i>r}u_iH_{ir}\right]^2 \pm 2u_r\sum_{i>r}u_iH_{ir},
    \end{align*}

    donde el $\pm$ será $+$ si $H_{rr}$ es positivo y $-$ si es negativo y de la misma forma en \eqref{eq4}. 
    
    Así, es claro que \eqref{eq3} y \eqref{eq4} son equivalentes. Y desde aquí es fácil ver como \eqref{eq1} se convierte en 
    \begin{equation*}
        f = \pm (V_1)^2 \pm\dots\pm (V_{r-1})^2 + (V_r)^2  \pm \sum_{i,j > r} V_iV_j\tilde{H}_{ij}(V_1,\dots,V_n)
    \end{equation*}
    para unos nuevos términos $\tilde{H}_{ij}$ simétricos en función del nuevo sistema de coordenadas.

    Esto concluye la inducción ya que tenemos un sumando cuadrático más y queda probada la existencia del sistema de coordenadas local buscado.
    
\end{proof}

\begin{remark}
    Durante la demostración queda implícita una parte del proceso. Al definir el nuevo sistema de coordenadas en el paso de inducción lo que estamos haciendo es definir una función que podemos llamar $l:\mathbb{R}^n \to \mathbb{R}^n$ que, dadas las coordenadas $(u_1(p),\dots, u_n(p))$ de un punto $p\in M$, las lleva a unas nuevas $(v_1(p),\dots, v_n(p))$ y es un difeomorfismo entre abiertos de $\mathbb{R}^n$. Así, si la carta que inducía el sistema de coordenadas $(u_1,\dots, u_n)$ era $\varphi : M \to \mathbb{R}^n$, la que induce $(v_1,\dots, v_n)$ será $l\circ \varphi$.
\end{remark}
\bigskip

El siguiente corolario será clave a la hora de demostrar uno de los teoremas principales de la teoría de Morse (ver \ref{cwvariedad}) y muestra una de las ventajas fundamentales de trabajar con puntos críticos no degenerados.

\begin{corollary}
        Sea $M$ una variedad diferenciable. Todo punto crítico no degenerado de una función suave de $M$ en $\mathbb{R}$ está aislado respecto de los demás puntos críticos.
\end{corollary}
\begin{proof}
    Sea $p\in M$ un punto crítico no degenerado de una función $f$. Existe un sistema de coordenadas local $(y_1,\dots, y_n)$ en un entorno $\mathcal{U}$ de $p$ tal que $y_i(p)=0,~~\forall i$ y 
    \begin{equation*}
        f=f(p)-(y_i)^2-\dots-(y_\lambda)^2 + (y_{\lambda +1})^2+\dots + (y_n)^2
    \end{equation*}
    Por el Lema de Morse (\ref{lemamorse}).\\

    Supongamos que existe otro punto crítico $q\in\mathcal{U}$, entonces
    \begin{equation*}
        \frac{\partial f}{\partial x_i}(q)=\pm 2q_i= 0\iff q_i= y_i(q)=0\iff q_i=p_i
    \end{equation*}
    Luego $p$ es el único punto crítico de f en $\mathcal{U}$.
    
\end{proof}

A continuación presentamos algunos puntos críticos degenerados y no degenerados junto con las gráficas de sus funciones correspondientes. Obsérvese que mientras los puntos no degenerados han de ser aislados, una función puede tener puntos críticos degenerados tanto aislados como no aislados.

\begin{example}
    \hfil
    \begin{enumerate}
        \item La función $f(x,y)=(x^2-1)^2 + (y^2-1)^2$ tiene nueve puntos críticos no degenerados. Es fácil comprobar que $\text{Hess}_pf$ es una matriz diagonal cuyos elementos no se anulan en ninguno de los puntos críticos. 
        \item La función $g(x,y)=x^3-3xy^2$ cuya gráfica se corresponde con la ``silla de mono'' posee un único punto crítico degenerado en el punto $(0,0)$. Basta ver que $\text{Hess}_{(0,0)}f$ tiene ceros en las cuatro entradas.
        \item Observamos que $h(x,y)= x^2y^2$ tiene infinitos puntos críticos degenerados, los de la forma $(0,y)$ y $(x,0)$.
    \end{enumerate}
\end{example}

\begin{figure}[H]
\begin{minipage}{0.35\textwidth}
    \centering
    \includegraphics[width=\textwidth]{/home/kike/Documents/Mates/TFG/introduction-to-morse-theory/memoria_libro/imagenes/grafica1.png}
    \caption{$f(x,y)=(x^2-1)^2 + (y^2-1)^2$}
\end{minipage}
\hfill
\begin{minipage}{0.35\textwidth}
    \centering
    \includegraphics[width=\textwidth]{/home/kike/Documents/Mates/TFG/introduction-to-morse-theory/memoria_libro/imagenes/grafica2.png}
    \caption{$g(x,y)=x^3-3xy^2$}
\end{minipage}

% --- Segunda fila ---
\centering
\begin{minipage}{0.35\textwidth}
    \includegraphics[width=\textwidth]{/home/kike/Documents/Mates/TFG/introduction-to-morse-theory/memoria_libro/imagenes/grafica3.png}
    \caption{$h(x,y)= x^2y^2$}
\end{minipage}
\end{figure}

Concluimos el capítulo definiendo los objetos que le dan nombre y de los que nos serviremos de ahora en adelante para describir la estructura de CW-complejo de cualquier variedad diferenciable.
\bigskip

\begin{definition}
 Diremos que una función suave $f: M\to \mathbb{R}$ es una \textit{función de Morse} si todos sus puntos críticos son no degenerados.
\end{definition}

\iffalse
{\color{blue} ¿Demostrar que para toda variedad de dimension finita existe una funcion de morse?}  
\fi